\section{Methodik}
% Struktur abgestimmt am 28.10.2026 (12 Seiten)

\subsection{Forschungsdesign}
\offen{4.1 -- experimentell und iterativ (v1--v4), mehrere Fälle mit derselben Methode; Überblick, welche Hypothese mit welchem Versuch geprüft wird}

\subsection{Auswahl der Anwendungsfälle}
% Satz 1 und Tabelle 4.1 abgesegnet am 28.10.2026
Der Ansatz wird an mehreren Anwendungsfällen untersucht, die sich in Branche, Rechtsgebieten und Zugänglichkeit der Vorschriften unterscheiden.

\begin{table}[H]
  \centering
  \small
  \renewcommand{\arraystretch}{1.3}
  \begin{tabularx}{\textwidth}{>{\bfseries\raggedright\arraybackslash}p{2.6cm} Y Y Y}
    \toprule
     & \textbf{Fall 1: Bewerbungsverfahren} & \textbf{Fall 2: Brandmeldetechnik} & \textbf{Fall 3: Vergleichsdaten Sai et al.} \\
    \midrule
    Organisation & öffentliche Hochschule in Bayern & Errichter bzw.\ Betreiber einer Brandmeldeanlage & Versicherung und Bank (Australien) \\
    Prozess & Besetzung einer befristeten Stelle mit KI-Vorauswahl & Planung bis Instandhaltung einer Brandmeldeanlage & Schadensbearbeitung; Kundenprüfung \\
    Rechtsgebiete & Datenschutz, KI/AI Act, Arbeits- und Dienstrecht & Bauordnungsrecht, technische Normen & Versicherungs- und Finanzaufsichtsrecht \\
    Zugang zu den Vorschriften & frei verfügbar & teilweise urheberrechtlich geschützt (DIN/VDE) & durch die Studie vorgegeben \\
    Rolle in der Arbeit & Hauptfall & explorativer Fall: Übertragbarkeit & Vergleichswert \\
    Fachliche Prüfung & Fachleute der HM (Datenschutz; KI und Arbeitsrecht offen) & Lehrbeauftragter für Brandschutz der HM, eigene Fachkenntnis & Fachleute der Studie \\
    Stand & umgesetzt & geplant & Verfügbarkeit prüfen \\
    \bottomrule
  \end{tabularx}
  \caption{Anwendungsfälle der Arbeit}
  \label{tab:anwendungsfaelle}
\end{table}
% Sätze 2-6 abgesegnet am 28.10.2026 ("Goldstandard" auf Wunsch durch "Referenzzuordnung" ersetzt)
Die Fälle wurden nach vier Kriterien ausgewählt.
Erstens sollen sie unterschiedliche Rechtsgebiete und Rechtsebenen abdecken, vom EU-Recht über Bundes- und Landesrecht bis zu technischen Normen.
Zweitens soll mindestens ein Fall Vorschriften enthalten, die nicht frei zugänglich sind.
Drittens müssen Fachleute erreichbar sein, die die \fb{referenzzuordnung} prüfen.
Viertens muss sich der Prozess in die drei Ebenen Prozess, Teilprozess und Task gliedern lassen.
% Sätze 7-9 abgesegnet am 30.10.2026
Die Fälle 1 und 2 wurden für diese Arbeit selbst entworfen.
Sie orientieren sich an typischen Abläufen in solchen Organisationen, geben aber keinen tatsächlichen Prozess einer bestimmten Organisation wieder.
Fall 3 übernimmt dagegen die Daten aus der Studie von Sai et al.~\cite{sai2024}.
\offen{Zeile \enquote{Stand} in Tabelle~\ref{tab:anwendungsfaelle} vor Abgabe aktualisieren oder entfernen}

\subsection{Fall 1: Bewerbungsverfahren an einer Hochschule}\label{sec:fall1}
% Sätze 1-3 abgesegnet am 28.10.2026
Der erste Anwendungsfall ist die Besetzung einer Stelle für eine wissenschaftliche Mitarbeiterin oder einen wissenschaftlichen Mitarbeiter an der Hochschule München.
Die Stelle ist auf drei Jahre befristet und dient der Promotion.
Für den Fall wird angenommen, dass die Hochschule ein KI-Werkzeug einsetzt, das eingehende Bewerbungen auswertet und in eine Rangfolge bringt.
\offen{Warum der Fall für die Arbeit sinnvoll ist: was er zum Test der Methodik beiträgt}
\offen{Hinweis: Eingaben und Referenz von Fall~1 sind in dieser Form zur fachlichen Prüfung vorgesehen (Abschnitt~\ref{sec:fachliche_pruefung})}
\subsubsection{Eingabe}\label{sec:fall1_eingabe}
% Sätze 4-8 und 10 abgesegnet am 18.11.2026 (Satz 9 gestrichen)
Der Agent erhält drei Eingaben, ein Profil der Organisation, den Prozess und die Annahmen zum Fall.
Das Profil beschreibt die Organisation, in der der Prozess abläuft.
Dazu gehören Sitz und Bundesland, Rechtsform, Größe und die Gremien der Mitbestimmung.
Die Hochschule München ist eine Körperschaft des öffentlichen Rechts in Bayern mit mehr als 1.000 Beschäftigten.
Sie hat einen Personalrat, eine Gleichstellungsbeauftragte und eine Schwerbehindertenvertretung.
Tabelle~\ref{tab:h1_profil} zeigt alle Angaben des Profils, die der Agent erhält (Datei \repodatei{h1/input/firmenprofil.json}).
\offen{Profil um Schwerbehindertenvertretung und Rolle beim KI-Tool (Betreiberin) ergänzen, Vorschlag in ERGEBNIS.md Abschnitt~8; Tabelle danach neu erzeugen. Entscheiden, ob die Tabelle hier oder im Anhang steht}
\offen{Prozess mit Teilprozessen und Tasks}
\offen{Annahmen (KI-Tool der Hochschule, kein Dienstleister außerhalb der EU, kein Talentpool)}

\subsubsection{Erwartung}\label{sec:fall1_erwartung}
\offen{Referenzzuordnung (Anhang~\ref{anh:vorschriften}), bewusst entfernte Schritte, Kontrollschritt}

\subsubsection{Abdeckung}\label{sec:fall1_abdeckung}
\offen{Tabelle der Rechtsebenen aus Abbildung~2.2, die der Fall abdeckt (Verweis auf Tabelle~\ref{tab:regelwerke} und Abschnitt~2.3)}

\subsection{Fall 2: Brandmeldetechnik}\label{sec:fall2}
\offen{Warum der Fall für die Arbeit sinnvoll ist: was er zum Test der Methodik beiträgt}
\offen{Entscheidung offen: Bekommt Fall~2 einen KI-Schritt (z.\,B. KI zur Bewertung von Meldungen und Täuschungsalarmen)? Sonst ist der AI Act in Fall~2 nicht anwendbar; Titel spricht von KI-gestützten Prozessen}
\offen{Normtexte sind geschützt, der Agent kann sie nicht laden, Prüfung von Hand über die HM-Bibliothek}
\subsubsection{Eingabe}\label{sec:fall2_eingabe}
\offen{Organisationsprofil, Prozess mit Teilprozessen und Tasks, Annahmen}

\subsubsection{Erwartung}\label{sec:fall2_erwartung}
\offen{Referenzzuordnung (Anhang~\ref{anh:vorschriften}), bewusst entfernte Schritte, Kontrollschritt}

\subsubsection{Abdeckung}\label{sec:fall2_abdeckung}
\offen{Tabelle der Rechtsebenen aus Abbildung~2.2, die der Fall abdeckt (Verweis auf Tabelle~\ref{tab:regelwerke} und Abschnitt~2.3)}

\subsection{Fall 3: Vergleichsdaten von Sai et al.}\label{sec:fall3}
\offen{Warum der Fall für die Arbeit sinnvoll ist: was er zum Test der Methodik beiträgt}
\offen{Daten: Versicherung (7 Teilprozesse, 31 Tasks, 342 veröffentlichte Absätze) und Bank (7 Teilprozesse, 19 Tasks, 311 Absätze), MIT-Lizenz, im Repo unter h1/input/sai. Ohne bewusst entfernte Schritte, daher nur für H1 und H2}
\subsubsection{Eingabe}\label{sec:fall3_eingabe}
\offen{Organisationsprofil, Prozess mit Teilprozessen und Tasks, Annahmen}

\subsubsection{Erwartung}\label{sec:fall3_erwartung}
\offen{Referenzzuordnung (Anhang~\ref{anh:vorschriften})}

\subsubsection{Abdeckung}\label{sec:fall3_abdeckung}
\offen{Tabelle der Rechtsebenen aus Abbildung~2.2, die der Fall abdeckt (Verweis auf Tabelle~\ref{tab:regelwerke} und Abschnitt~2.3)}

\subsection{Fachliche Prüfung}\label{sec:fachliche_pruefung}
% Gliederung 4.6-4.11 abgestimmt am 18.11.2026
\offen{Wer prüft je Fall (Rolle, Fachgebiet), Ablauf, Prüfbogen (Anhang~D), Ergebnis}
\offen{Fachliche Prüfung der Referenzzuordnung (Fall~1, Fall~2) durch [Fachperson, Rechtsgebiet] ausstehend. Fall~3 wurde bei Sai et al.\ geprüft. Grundlage für H1, H2 und H3}
\offen{Fachliche Prüfung der Soll-Entscheidungen der Referenzfälle für die Prüfregeln ausstehend. Grundlage für H5, H6 und H7}

\subsection{Agent}\label{sec:agent}
\offen{Modell, Parameter, Ausgabeformat}
\offen{Stufe~1 (Regelwerke) und Stufe~2 (Zuordnung je Task, Lückenmeldungen)}
\offen{Werkzeuge und amtliche Quellen; Prompts in Anhang~C}

\subsection{Durchführung}\label{sec:durchfuehrung}
\offen{Konfiguration A (Agent bestimmt Regelwerke selbst) und B (Regelwerke vorgegeben)}
\offen{Versionen v1--v4 mit ihren Änderungen; 5 Läufe je Version und Konfiguration}
\offen{Docker, Protokoll, Repository, Budgetgrenze und Kosten}
\offen{Satz: Verweis auf Anhang~\ref{anh:testaufbau} (Eingaben und Referenz je Hypothese)}

\subsection{Auswertung H1 bis H3}\label{sec:auswertung_h1_h3}
\subsubsection{Bewertung je Fund}
\offen{Kategorien: korrekt, neutral, falscher Absatz, falscher Task, nicht in der Referenz; Prüfung auf erfundene Normen}
\subsubsection{H1: Regelwerke}
\offen{Trefferquote und Genauigkeit der Regelwerke, Schwelle 0,9}
\offen{Satz: Herleitung der Referenz für H1 aus der Referenzzuordnung, Verweis auf Tabelle~\ref{tab:h1_referenz}}
\subsubsection{H2: Zuordnung}
\offen{Trefferquote auf Task-Ebene, Teilprozess- und Prozessebene ergänzend; Schwelle offen (0,62 oder 0,69)}
\subsubsection{H3: Fehlende Schritte}
\offen{Erkannte Lücken, Fehlmeldungen, Kontrollschritt; Grenze für Fehlmeldungen offen}

\subsection{Aufbau und Auswertung H4}\label{sec:h4}
\offen{Eingabe ohne Prozess, Referenz aus Reihenfolge-Paaren \enquote{A vor B} je Fall, bedingter und strenger Wert, Schwelle 80\,\%}

\subsection{Aufbau und Auswertung H5 bis H7}\label{sec:h5_h7}
\offen{Referenzfälle mit Soll-Entscheidung (gemeinsam für H5 bis H7)}
\subsubsection{H5: Prüfregeln}
\offen{Regeln aus den Zuordnungen erzeugen, Übereinstimmung mit der Soll-Entscheidung}
\subsubsection{H6: Laufzeitmonitor}
\offen{KI-Agent führt den Prozess aus, Monitor prüft jede Aktion; blockierte unzulässige und fälschlich blockierte zulässige Aktionen}
\subsubsection{H7: Gesamtkette}
\offen{Ablauf wie H6 mit den ungeprüften Ergebnissen aus H2 und H3, Vergleich mit H6}
